% \maketitle
\documentclass[lettersize,journal]{IEEEtran}
\usepackage{amsmath,amsfonts,amssymb}
\usepackage{algorithmic}
\usepackage{algorithm}
\usepackage{array}
\usepackage{textcomp}
\usepackage{stfloats}
\usepackage{url}
\usepackage{verbatim}
\usepackage{graphicx}
\usepackage{booktabs}
\usepackage{multirow}
\usepackage{dashrule}

\usepackage{subcaption} 
\usepackage{cite}
\usepackage{etoolbox} 

% =========================================================================

% =========================================================================
\usepackage{orcidlink}
\usepackage{hyperref}
\hypersetup{
    colorlinks=true,
    linkcolor=blue, 
    citecolor=blue,  
    urlcolor=black   
}


\makeatletter

\renewcommand\citeform[1]{\textcolor{blue}{#1}}
\renewcommand\citepunct{\textcolor{blue}{, }}
\renewcommand\citedash{\textcolor{blue}{--}}


\def\citeleft{\textcolor{blue}{[}}
\def\citeright{\textcolor{blue}{]}}
\makeatother


\hyphenation{op-tical net-works semi-conduc-tor IEEE-Xplore}

\begin{document}

\title{QAG-LDM: Quality-Aware Gated Multi-source Latent Diffusion Model for Music Generation and Source Extraction}

\author{Qinran~Lin\,\orcidlink{0000-0000-0000-0000},
        Yiqun~Wang\,\orcidlink{0003-1942-5597},~\IEEEmembership{Member,~IEEE},
        Hao~Zhang\,\orcidlink{0009-0005-4966-8081},
        and~Yong~Li\,\orcidlink{0000-0003-1584-1622},~\IEEEmembership{Member,~IEEE}% <- 加上了逗号，最后用%防止多余空格

        % <-this % stops a space

% \thanks{Manuscript received XXX XX, 20XX; revised XXX XX, 20XX; accepted XXX XX, 20XX. Date of publication XXX XX, 20XX; date of current version XXX XX, 20XX.The associate editor coordinating the review of this article and approving it for publication was XXX.(Corresponding author: Yong Li.)}
% \thanks{Qinran~Lin, Yiqun~Wang, Hao~Zhang, and Yong~Li are with the School of Computer Science and Technology, Chongqing University, Chongqing 400044, China (e-mail: 202414021050t@stu.cqu.edu.cn; yiqun.wang@cqu.edu.cn; 20241401018@stu.cqu.edu.cn; yongli@cqu.edu.cn).}%
\thanks{Corresponding author: Yong Li.}
\thanks{Qinran~Lin, Yiqun~Wang, Hao~Zhang, and Yong~Li are with the School of Computer Science and Technology, Chongqing University, Chongqing 400044, China (e-mail: 202414021050t@stu.cqu.edu.cn; yiqun.wang@cqu.edu.cn; 20241401018@stu.cqu.edu.cn; yongli@cqu.edu.cn).}%
\thanks{Digital Object Identifier }
}
\markboth{IEEE TRANSACTIONS ONAUDIO,SPEECHANDLANGUAGEPROCESSING,VOL.XX,20XX}%
{Shell \MakeLowercase{\textit{et al.}}: A Sample Article Using IEEEtran.cls for IEEE Journals}

\maketitle
\begin{abstract}
Music generation has attracted increasing attention in recent years, driven by the rising demand for high-quality cultural and creative content in the context of rapid advancements in artificial intelligence. Current music generation models, especially the multi-source diffusion model (MSDM) series, have integrated the processes of total generation, partial generation, and source extraction. However, the prevalence of quality-imbalanced waveforms in multi-stem datasets hinders models from learning robust generative priors, leading to instability. To address this gap, we propose QAG-LDM, a novel framework that enhances multi-stem music generation and source extraction with quality-aware conditioning and adaptive gated modulation. Specifically, quality scores based on a frozen COCOLA module are converted into learnable quality tokens through our proposed QT-Encoder, which are prepended to the latent sequence alongside audio and timestep embeddings. 
% The gating mechanism is applied to the output of self-attention layers, addressing the attention sink phenomenon on the timestep token we identified in experiments. 
We further identify an attention sink phenomenon centered on the timestep token and introduce a gated self-attention modulation strategy to adaptively regulate attention propagation across layers.
Specifically, gates modulate the attention output, effectively preventing excessive concentration of attention on initial tokens without weakening the original intended meaning. 
% To ensure effective quality guidance, we incorporate an auxiliary quality score predictor during training, enabling the model to learn robust quality-conditioned generation. 
To facilitate effective quality-aware generation, we introduce an auxiliary quality prediction objective that explicitly guides the model toward robust quality-aware representation learning.
Through extensive ablation studies on the Slakh2100, Musdb18, and MoisesDB datasets, our method achieves superior performance compared to all baseline variants on FAD, Log-Mel L1, and COCOLA scores across standard MSDM tasks.
\end{abstract}

\begin{IEEEkeywords}
Multi-source Music Generation, Diffusion Transformer, Quality-Aware Synthesis, Gated Attention, Controllable Audio Generation.
\end{IEEEkeywords}

% % 数据质量不一且模型不分好坏（第二段） 
% 他人尝试用规模或结构补救（第三段）
% 但 MSDM 架构仍无视质量感知，导致分布不理想（第四段） 
% 我们提出质量感知的解耦控制（方案）。

\section{Introduction}
\IEEEPARstart{T}{he} main goal of music generation is to transform abstract textual information into audio music, which provides new tools for art creation. Early works \cite{Oord_2016,engel2019gansynth} concentrated on using various generative models to represent music as quantized waveforms or spectral features. The approach MSDM\cite{mariani2023multi} for the first time enables the total generation, source extraction, and partial generation tasks within a single framework, improving interactive methods for generating music. 

High-quality music data is limited by copyright and production expenses. Professional music production necessitates costly hardware, high-sampling-rate recording equipment, and musicians possessing extensive artistic skill. Consequently, this reliance on specialized resources results in a significantly smaller volume of high-quality annotated data accessible to open-source communities compared to the image or text domains. More critically, existing large-scale music datasets often exhibit high heterogeneity in terms of audio quality, encompassing a mixture of studio-grade recordings and low-fidelity samples. The lack of quality-aware supervision prevents models from learning a clear mapping of high-quality distributions, especially when operating under constrained data environments where noisy or low-quality samples can disproportionately degrade the generative performance.

Current research efforts in music generation and source separation aim to enhance the performance of the model in the aforementioned constrained data environment. Specifically,  AudioSep \cite{liu2024separate} mitigates the limitations of small-scale supervised data by utilizing multimodal pre-training on massive datasets, establishing a foundation model capable of separating arbitrary sound sources through natural language queries. Research \cite{plaja2025leveraging} utilizes a regressive prior to guide the diffusion process for separation in specialized repertoires. Other efforts, such as SoloAudio \cite{wang2025soloaudio}, leverage synthetic audio and Transformer-based architectures to enhance zero-shot extraction capabilities, while FLOSS \cite{scheibler2025source} utilizes flow matching with strict mixture consistency constraints to solve the ill-posed separation problem via equivariant generative modeling. 
Beyond source separation, generative tasks also face challenges in constrained data environments, particularly for style-specific and multi-task content creation. As a representative work in efficient generation, MusicFlow \cite{prajwal2024musicflow} employs cascaded flow matching to achieve high-quality synthesis with reduced parameters. Addressing the trade-off between stem flexibility and inference speed, STEMPHONIC \cite{wu2026stemphonic} introduces a shared noise latent mechanism that enables the one-pass generation of a variable set of synchronized music stems. Furthermore, research \cite{cao2025ai} has explored using Diffusion Transformer (DiT) to generate traditional Chinese music from videos to overcome the scarcity of culturally aligned data, while \cite{yu2025improving} efforts have introduced Gaussian priors to calibrate inference-time optimization for vocal style transfer in limited data regimes.

Within current MSDM series frameworks, data are organized in different formats, such as the standard four-stem format \cite{mariani2023multi} and the triplet format \cite{chae2025mge}. However, these frameworks primarily focus on architectural integration for multi-source tasks while neglecting the data quality challenges inherent in music production mentioned above. Specifically, the training paradigm of existing MSDM models remains essentially "quality-agnostic," implicitly assuming a uniform high-fidelity distribution across all training samples. When researchers introduce new modeling approaches, they often prioritize architectural feasibility and structural alignment while largely overlooking the intrinsic quality relationships within music organization formats. These models frequently lack the quality-aware mechanism to perceive or prioritize the quality information embedded within these structured combinations, a deficiency that ultimately hinders their ability to align with human auditory perceptions, such as coherence. For example, when the data is organized in triplet format (see Figure \ref{fig:cocola_distribution}), the upper and lower limits of the data's performance on auditory perceptions vary widely with an unsatisfactory distribution, proving that a lack of quality awareness leads to an uncontrollable blending of high- and low-fidelity features.
 
\begin{figure}[t]
    \centering
    % 如果需要裁剪，取消注释下行并调整数值：trim = left bottom right top
    \includegraphics[trim=0bp 0bp 0bp 25bp, clip, width=\linewidth]{Figure_1.png}
    % \includegraphics[width=\linewidth]{Figure_1.png}
    \caption{COCOLA score distributions across the Slakh2100 (blue), MUSDB18 (orange), and MoisesDB (green) datasets. For each audio track, one stem is randomly selected as the target source, with the remainder aggregated as the sub-mixture. To unify stem semantics for joint dataset training, specific track mappings are applied: "Strings (continued)" is mapped to "Strings" in Slakh2100, and "other key" is mapped to "synths" in MoisesDB.}
    \label{fig:cocola_distribution}
\end{figure}

To address these gaps, we propose QAG-LDM, a Quality-Aware Gated Multi-source Latent Diffusion Model for Music Generation and Source Extraction. 
% Specifically, text embeddings are processed via a frozen CLAP encoder, while quality scores $q \in [0,1]$ are converted into learnable quality tokens through a QT-Encoder, enabling the model to perceive quality information. 
Specifically, we formulate quality-aware conditioning by projecting continuous quality scores $q \in [0,1]$ into learnable quality tokens via a dedicated QT-Encoder, allowing the diffusion model to explicitly model perceptual quality information.
These tokens are then processed via dedicated attention branches, where a gated attention mechanism is applied to the output of self-attention layers to adaptively modulate the influence of quality conditioning and mitigate the attention sink phenomenon on timestep token observed in baseline models. 
% Furthermore, an auxiliary loss is applied to ensure effective quality guidance, while the text-conditioned CFG dropout is inherited from the baseline. 
Furthermore, an auxiliary quality supervision objective is introduced to strengthen quality-aware guidance while preserving the text-conditioned CFG dropout strategy from the baseline framework.
This training paradigm is dataset-agnostic.

We evaluated the model's performance within our framework through a series of controlled experiments. The comparative testing outcomes indicate that our model attains a low FAD score across the entire generation, closely mirroring the distribution of real-world audio. Furthermore, during source extraction, the model achieved high accuracy in spectrum reconstruction when extracting target stems from mixed audio. The audio track generated under specific parameters showed the highest semantic consistency compared to the other background tracks in the partial generation.
Ablation experiments further validated the efficacy of our design, demonstrating that our method achieves the best performance among all tested approaches. The quality token mechanism plays a primary role in enabling quality-aware generation, while the gated attention mechanism optimizes the model's expressive capacity. Our contributions are as follows:

\begin{itemize}
    \item We propose a quality-aware paradigm using our designed QT-Encoder that maps normalized COCOLA scores into high-dimensional latent tokens, enabling models to train via the dataset's quality-guided process, improving musicality and audio quality in music generation and extraction tasks.
    
    \item We propose a Gated Self-Attention (GSA) modulation strategy within the DiT backbone. By enabling the model to perceive quality information and text information combined with a gated mechanism, effectively mitigating the attention sink phenomenon on timestep tokens.
    
    \item We address the problem of ambiguous coordination among randomly combined mixed stems within a single track in large music datasets. Within the MSDM series, our method achieves SOTA performance in total generation and source extraction, while specifically attaining SOTA on the COCOLA score for partial generation.

\end{itemize}

\section{Related Work}
\subsection{Audio Generation Models }
Early research such as WaveNet \cite{Oord_2016} and SampleRNN \cite{Mehri_2016} focused on autoregressive models. Researchers directly model raw audio waveforms to capture temporal dependencies to some extent. Nevertheless, their scalability was constrained by the pointwise sampling method. Later studies \cite{Huang_2018,Dhariwal_2020} incorporate the Transformer architecture \cite{vaswani2017attention} to capture long-range musical dependencies. Others, by employing methods such as GANs and variational autoencoders (VAEs) \cite{kumar2019melgan,kong2020hifi}, enhanced previous music generation models. Although these researchers start to focus on timbre quality, their approaches had limitations in how they represented the overall structure of music, and they struggled to maintain stability over long audio. Since the work of \cite{Nichol_2021}, diffusion models have provided a more stable and expressive way to generate audio, as demonstrated in applications such as \cite{kong2020diffwave}. Despite this, directly using diffusion in the waveform space is still computationally expensive. As a result, LDM \cite{Rombach_2021} uses autoencoders, including SoundStream \cite{zeghidour2021soundstream} and EnCodec \cite{defossez2022high}, to compress audio into latent spaces, enabling more efficient generation. Building on LDM and CLAP's cross-modal alignment, the AudioLDM \cite{Liu_2023} series has become a mainstream method for modern music generation, balancing computational efficiency with semantic consistency.

\subsection{Multi-Track Music Audio Modeling }
Recent research has begun to consider music as a structured combination, including discrete representations. StemGen \cite{parker2024stemgen} employs an iterative non-autoregressive Transformer for text-guided generation. MusicGen-Stem \cite{rouard2025musicgen} uses stem vector quantization and an autoregressive decoder to generate specific musical components, such as basslines and percussion. 
At the same time, researchers began to focus on diffusion for modeling. Furthermore, subsequent research, such as MusicLDM \cite{chen2024musicldm} and Tango 2 \cite{majumder2024tango}, has investigated cross-modal generation, with Tango 2 specifically focusing on audio-video generation. Jen-1 Composer \cite{yao2025jen} uses latent diffusion on standard stems to create unified multi-stem music. A foundational work in this area was MSDM \cite{mariani2023multi}. This unified framework uses additive mixing and sampling based on ordinary differential equations. Later models have built upon this original work. GMSDI \cite{postolache2024generalized} uses variable stems and text-to-waveform inference. In contrast, MSD-LD \cite{karchkhadze2025simultaneous} adds Classifier-Free Guidance (CFG) \cite{ho2022classifier} to their work to achieve more controlled generation. Therefore, MGE-LDM \cite{chae2025mge} broadens its scope of application, moving from predefined to arbitrary categories, thereby creating a more flexible research environment.

\subsection{Conditional Control and Quality Enhancement }
The nearest music generation models are increasingly emphasizing conditional control. They allow users to prompt various conditions, including textural conditions, musical characteristics, instrument types, rhythmic patterns, and emotional states \cite{mitra2025music}.  WavCaps \cite {mei2024wavcaps} and Make-an-Audio 2 \cite {huang2023make} leverage ChatGPT \cite {achiam2023gpt} to produce more accurate textual descriptions, reinforcing cross-modal alignment. LP-MusicCaps \cite {doh2023lp} and Music-LLaMA \cite {liu2024music} have designed structured description systems for music, offering more detailed semantic expressions for conditional control. 
Issues such as low sampling rates or mono representations frequently arise during the audio generation process. Consequently, the synthesized audio often suffers from degraded quality. To address these, models like MusicGen and Stable Audio \cite{evans2025stable} directly model stereo audio at 32 kHz or 44.1 kHz. This design helps improve tonal detail, spatial perception, and overall perceived quality. To further handle quality-imbalanced data, QA-MDT \cite{li2024quality} introduces a quality-aware training paradigm and adapts the Masked Diffusion Transformer (MDT) to strengthen the spatial correlation of music signals. Several studies have attempted to evaluate generated audio from multiple perspectives to achieve better results. For generating online music accompaniment, the deep reinforcement learning algorithm RL-Duet \cite{jiang2020rl} requires real-time generation of coherent accompaniment that matches human input \cite{wu2025streaming}. During evaluation, metrics such as real-time response speed and interaction fluency are key focuses alongside generation quality. COCOLA \cite {ciranni2025cocola} designed a contrastive learning-based encoder specifically for objectively evaluating the harmonic and rhythmic coherence of music accompaniment generated by models. This addresses the limitation of existing metrics (such as FAD or sub-FAD) that prioritize global quality over track-to-track coherence. DUO-TOK \cite{lin2025duo} proposes a dual-track semantic music segmenter that uses a four-stage, source-aware encoder-decoder architecture to improve the quality and consistency of the generated dual-track music.

\section{Preliminary}
\subsection{Latent Diffusion Model Framework for Multi-Music Generation }
MSDM first proposed a method for generating waveform audio for standard stems (bass, drums, guitar, and piano) using diffusion models, with subsequent approaches extending this concept. The latent-based diffusion model first employs a pre-trained VAE \cite{kingma2019introduction} to compress the raw audio $x_0$ into a high-dimensional latent space, yielding the latent representation $z_0$. The two key processes of diffusion models, diffusion and denoising, operate efficiently within this latent space: During diffusion, Gaussian noise $\epsilon$ is incrementally added to $z_0$ at each timestep $\tau$, ultimately generating a noisy latent representation $z_{\tau}$. The core of the reverse denoising process involves training a denoising network $f_{\theta}$ (typically a DiT\cite{peebles2023scalable}) to estimate the noise or velocity target $v_{\tau}$, thereby recovering $z_0$ from $z_{\tau}$. When using a $v$-objective loss function, it is generally defined as $\mathcal{L}(\theta)=\mathbb{E}_{z_{0},\epsilon,\tau}\lVert f_{\theta}(z_{\tau},\tau,c)-v_{\tau}\lVert _{2}^{2}$, where $c$ represents the conditional embedding required by the model. During inference, the model generates target latent representations $\hat{z}_0$ from random noise via iterative sampling. These are input into VAE to obtain the final audio waveform $\hat{x}_0$. Some approaches incorporate conditional generation \cite{dhariwal2021diffusion}, designing the control signal $y$ that influences the denoising function $D_\theta$ to generate content $D_\theta(z_t, t, y)$. The subsequently proposed CFG enhances this process by introducing an unconditional model $D_\theta(z_t, t, \varnothing)$, yielding the following relationship:$\bar{D}_\theta(z_t, t, y) = D_\theta(z_t, t, \varnothing) + s(D_\theta(z_t, t, y) - D_\theta(z_t, t, \varnothing))$. Here, $\varnothing$ denotes the absence of a specific control signal and $s \geq 1$ represents the guidance scale. $y$ is replaced by $\varnothing$ with probability $p_{\mathit{uncond}}$ while training. MGE-LDM simultaneously learns and models the joint distribution $p_{\theta}(z^{(m)},z^{(u)},z^{(s)})$ of three interrelated latent representations in the latent space: the complete mixed audio $z^{(m)}$, the selected individual sound source $z^{(s)}$, and the residual portion $z^{(u)}$.

\subsection{Coherence-Oriented Contrastive Learning for Audio }
Coherence-Oriented Contrastive Learning for Audio (COCOLA) \cite {ciranni2025cocola} is a contrastive learning encoder designed to objectively measure the harmonic and rhythmic coherence of music accompaniment generated by models. This addresses the limitation of existing metrics (FAD/Sub-FAD), which prioritize overall quality over track-to-track coherence. The encoder is trained by maximizing the consistency (positive samples) between two non-overlapping track sub-mixes $m_1^k$ and $m_2^k$ within the same music window in the embedding space, while minimizing similarity (negative samples) across different windows. The COCOLA encoder $f_{\theta}$ is designed based on the EfficientNet-B0 \cite{tan2019efficientnet} architecture, which maps inputs to embedding vectors. It optionally operates on features after harmonic-percussion separation to enhance performance and distinguish between harmonic and rhythmic coherence factors. The COCOLA score measures the similarity between the target track embedding $h_y$ and the generated accompaniment embedding $h_{\tilde{x}}$, defined as $sim(h_{y},h_{\tilde{x}})=(h_{y})^{T}Wh_{\tilde{x}}$. This score shows a strong statistical correlation with MOS (the Pearson correlation coefficient is $r=0.54$).

\begin{figure*}[t]
    \centering
    % 左边：长图
    \begin{minipage}[c]{0.68\textwidth} % 稍微缩小一点宽度给虚线留空间
        \centering
        \includegraphics[trim=50bp 20bp 80bp 60bp, clip, width=\textwidth]{model -V3.pdf}
    \end{minipage}
    \hfill
    % --- 插入虚线 ---
    \begin{minipage}[c]{0.01\textwidth}
        \centering
        \rotatebox{90}{\hdashrule{7cm}{0.5pt}{2pt 2pt}} % 6cm是高度，0.5pt是粗细，2pt 2pt是虚线间距
    \end{minipage}
    \hfill
    % --- 插入虚线结束 ---
    % 右边：宽图
    \begin{minipage}[c]{0.28\textwidth}
        \centering
        \includegraphics[trim=300bp 30bp 300bp 5bp, clip, width=\textwidth]{GA.pdf}
    \end{minipage}
    \caption{Overview of the proposed framework. \textbf{Left}: Total pipeline with QT-Encoder as detailed in Section~\ref{4.2}. \textbf{Right}: Detailed GSA mechanism as detailed in Section~\ref{4.3}. The gated output is added via a residual connection and passed through a feed-forward network, producing the updated latent representation.} 
    \label{fig2}
\end{figure*}

\section{Method}

\subsection{Quality-Aware Conditioning Definition}
Due to the inherent structural variance in multi-stem audio data, naive concatenation or random mixing of source tracks can introduce semantic misalignment. Such inconsistency impairs the model's ability to learn stable relationships among stems. To provide an explicit and learnable quality signal, we formalize an explicit quality-aware mechanism.

Formally, we define the training domain as $\mathcal{D}=\{\mathit{m}_i\}$, where each joint sample $\mathit{m}_i$ is parameterized as a structural triplet:
  \[
    \mathit{m}_i=\{x_i^{(m)},x_i^{(u)},x_i^{(s)}\},
  \]
corresponding to the mixture, sub-mixture, and source stems, respectively. For each track $i$, let
  $\{x_j\}_{j\in\mathcal{J}}$ 
denote the collection of time-domain stems, where the cardinality of $\mathcal{J}$ may vary depending on the number of instruments. The mixture waveform is defined as
  \begin{equation}
    x_i^{(m)}=\sum_{j\in\mathcal{J}}x_j.
  \end{equation}
Given a selected source subset $\mathcal{K}\subseteq\mathcal{J}$, we define the source and sub-mixture signals as
  \begin{equation}
    x_i^{(s)}=\sum_{k\in\mathcal{K}}x_k,\qquad
    x_i^{(u)}=x_i^{(m)}-x_i^{(s)}.
  \end{equation}

To map raw acoustic waveforms into a computationally tractable latent manifold, each signal in the structural triplet is projected via a pretrained VAE encoder $E$ to derive continuous latent representations:
  \begin{equation}
    z_i^{(v)}=E(x_i^{(v)}),\qquad v\in\{m,u,s\},
  \end{equation}
where $z_i^{(v)}\in\mathbb{R}^{C\times L}$, with $C$ and $L$ denoting the latent channel dimension and temporal length, respectively.

To quantify the structural compatibility between the sub-mixture and source stems, we evaluate the pair
  $(x_i^{(u)},x_i^{(s)})$ 
using a frozen COCOLA module~\cite{ciranni2025cocola}, producing a scalar quality score
  \begin{equation}
    q_i=f_{\mathrm{COCOLA}}(x_i^{(u)},x_i^{(s)}).
  \end{equation}
To eliminate dataset-specific scale discrepancies and establish a unified conditional distribution during joint dataset training, the quality scalar is bounded through a dataset-level min-max clipping operation:
  \begin{equation}
    \tilde{q}_i=
    \mathrm{clip}\!\left(
    \frac{q_i-q_{\min}}{q_{\max}-q_{\min}},\,0,\,1
    \right),
  \end{equation}
where $q_{\min}$ and $q_{\max}$ denote the minimum and maximum quality scores estimated across all combined datasets. The resulting $\tilde{q}_i\in[0,1]$ operates as an explicit, continuous quality prior for each training sample, establishing the mathematical foundation for the subsequent quality-aware joint regulation paradigm.

% Then we sample a noise level $\tau \sim \mathcal{U}([\tau_{\min},1])$ for training, and construct a per-track timestep vector $\boldsymbol{\tau}\in\mathbb{R}^{3}$ in a probabilistic manner to unify total generation, partial generation, and source extraction \cite{chae2025mge}:
Specifically, given a global noise level $\tau \sim \mathcal{U}([\tau_{\min},1])$ to unify the formulation of total generation, partial generation, and source extraction, a joint timestep vector $\boldsymbol{\tau}\in\mathbb{R}^{3}$  is defined probabilistically as \cite{chae2025mge}:
  \begin{equation}
    \boldsymbol{\tau}=
    \begin{cases}
    (\tau,\tau,\tau), & \text{with probability } 1-p_{\mathrm{drop}},\\
    (0,\tau,\tau), & \text{with probability } p_{\mathrm{drop}}/3,\\
    (\tau,0,\tau), & \text{with probability } p_{\mathrm{drop}}/3,\\
    (\tau,\tau,0), & \text{with probability } p_{\mathrm{drop}}/3,
  \end{cases}
  \end{equation}
where $p_{\mathrm{drop}}$ denotes the probability of randomly selecting one observed track by setting its timestep to zero.

% The concatenated latent triplet $\{z_i^{(m)}, z_i^{(u)}, z_i^{(s)}\} \in \mathbb{R}^{3C \times L}$ is transposed and projected to the transformer hidden dimension $D$ through a linear layer, yielding the latent token sequence $Z_i \in
%   \mathbb{R}^{N \times D}$. The timestep vector $\boldsymbol{\tau}$ 
% is first mapped to a learnable timestep token
%   \begin{equation}
%     T_i \in \mathbb{R}^{1\times D}.
%   \end{equation}
The global input feature space is subsequently constructed by channel-wise concatenation and linear dimension adaptation. The multi-track latent representation $\{z_i^{(m)}, z_i^{(u)}, z_i^{(s)}\} \in \mathbb{R}^{3C \times L}$ is mapped to the transformer hidden dimension $D$, defining the structural audio token sequence $Z_i \in \mathbb{R}^{N \times D}$. Concurrently, the joint multi-task timestep vector $\boldsymbol{\tau}$ is parameterized into a continuous temporal condition token: \begin{equation}
T_i \in \mathbb{R}^{1\times D}.
\end{equation}

\subsection{Quality Tokenization via QT-Encoder and Joint Diffusion Training}
\label{4.2}
Building upon the normalized quality score $\tilde{q}_i$, we introduce the Quality Token Encoder (QT-Encoder) to transform this scalar signal into a sequence of prefix tokens that are directly injected into the diffusion transformer as shown in Figure~ \ref{fig2}. The normalized quality score is transformed into quality tokens through sinusoidal encoding followed by a learnable multilayer perceptron. Specifically, we define
  \begin{equation}
  \begin{aligned}
    \Gamma(\tilde{q}_i) = \big[ & \sin(2\pi 2^0 \tilde{q}_i), \cos(2\pi 2^0 \tilde{q}_i), \ldots, \\
    & \sin(2\pi 2^{F-1}\tilde{q}_i), \cos(2\pi 2^{F-1}\tilde{q}_i) \big] \in \mathbb{R}^{2F},
  \end{aligned}
  \end{equation}
and map it to quality tokens
  \begin{equation}
    e_q=\gamma\,MLP_{\theta}(\Gamma(\tilde{q}_i)),
  \end{equation}
where $\gamma\in\mathbb{R}$ is a learnable scalar that controls the strength of the injected quality signal. Since the MLP directly maps to the joint multi-track hidden dimension $D$, the output is reshaped as $Q_i\in\mathbb{R}^{M\times D}$. The timestep token and quality tokens are prepended to the latent token sequence before entering the transformer:
  \begin{equation}
    X_i=[\,T_i;\,Q_i;\,Z_i\,] \in\mathbb{R}^{(1+M+N)\times D}.
  \end{equation}

During training, for a random subset of samples, the quality token $Q_i$ are replaced by a learnable mask token with probability $p_m$. This masked-quality strategy encourages the model to preserve quality-related information in its hidden representations rather than relying solely on direct token injection. To supervise this behavior, we attach an auxiliary quality prediction head to the final hidden states corresponding to the quality-token positions. Let
  $\bar{h}_i^{(q)}$ 
denote the mean-pooled hidden representation over the quality-token positions in the final transformer layer. The predicted quality score is defined as
  \begin{equation}
  \hat{q}_i=\sigma\bigl(g(\bar{h}_i^{(q)})\bigr),
  \end{equation}
where $g(\cdot)$ is a lightweight prediction head and $\sigma(\cdot)$ is the sigmoid function. The auxiliary regression loss is computed only for samples whose quality tokens are masked:
  \begin{equation}
    \mathcal{L}_{\mathrm{q}}=
    \frac{1}{|\mathcal{B}_{\mathrm{mask}}|}
    \sum_{i\in\mathcal{B}_{\mathrm{mask}}}
    \left\|\hat{q}_i-\tilde{q}_i\right\|_2^2,
  \end{equation}
where $\mathcal{B}_{\mathrm{mask}}$ denotes the subset of masked samples in the training batch.

We employ velocity prediction ($v$-prediction) as the diffusion objective \cite{salimans2022progressive}. Given the clean latent representation $z_0$ and per-track timestep vector $\boldsymbol{\tau}$, Gaussian noise is injected into the latent space via cosine scheduling:
  \begin{equation}
    z_{\boldsymbol{\tau}}=\alpha_{\boldsymbol{\tau}}\odot z_0+\beta_{\boldsymbol{\tau}}\odot\epsilon,
    \qquad
    \epsilon\sim\mathcal{N}(0,I),
  \end{equation}
where $\odot$ denotes element-wise multiplication, and the coefficients are applied track-wise according to
  \begin{equation}
    \alpha_{\boldsymbol{\tau}}=\cos\!\left(\frac{\pi}{2}\boldsymbol{\tau}\right),
    \qquad
    \beta_{\boldsymbol{\tau}}=\sin\!\left(\frac{\pi}{2}\boldsymbol{\tau}\right).
  \end{equation}
Under this noise parameterization, the corresponding target velocity is
  \begin{equation}
    v_{\boldsymbol{\tau}}
    =
    \alpha_{\boldsymbol{\tau}}\odot\epsilon
    -
    \beta_{\boldsymbol{\tau}}\odot z_0.
  \end{equation}
Because one track may be explicitly observed by setting its timestep to zero, we apply the reconstruction loss only to the noisy tracks. The diffusion loss is therefore defined as
  \begin{equation}
    \mathcal{L}_{\mathrm{diff}}
    =
    \mathbb{E}_{z_0,\epsilon,\boldsymbol{\tau}}
    \left\|
    f_{\theta}(z_{\boldsymbol{\tau}},\boldsymbol{\tau},C_i,Q_i)
    -
    v_{\boldsymbol{\tau}}
    \right\|_2^2.
  \end{equation}
The final training objective combines the diffusion objective with the auxiliary quality supervision:
  \begin{equation}
    \mathcal{L}
    =
    \mathcal{L}_{\mathrm{diff}}
    +
    \lambda_{\mathrm{q}}\mathcal{L}_{\mathrm{q}},
  \end{equation}
where $\lambda_{\mathrm{q}}$ controls the contribution of the quality supervision term.

% \subsection{Gated Self-Attention Modulation in DiT}
% \label{4.3}
% Our denoising backbone is a latent diffusion transformer \cite{evans2024long} that operates on the concatenated three-track latent sequence $X_i$. The self-attention operation is defined as
%   \begin{equation}
%     Attn(Q,K,V)=\mathrm{softmax}\!\left(\frac{QK^\top}{\sqrt{d_k}}\right)V,
%   \end{equation}
% where $d_k$ is the head dimension. In multi-head form, the self-attention output is
%   \begin{equation}
%     SA(X)=
%     \mathrm{Concat}\!\left(head_1,\ldots,head_h\right)W^{O},
%   \end{equation}
% with
%   \begin{equation}
%     head_i=
%     Attn(XW_i^Q,XW_i^K,XW_i^V).
%   \end{equation}

% To further modulate attention, especially the attention sink phenomenon on timestep token, we employ a query-dependent output gate. Given the attention input $X$, a learnable gating function produces an element-wise gate
%   \begin{equation}
%     G(X)=\sigma\!\left(XW^{G}\right),
%   \end{equation}
% where $W^{G}$ refers to the learnable parameters of gate and $\sigma$ is an activation function (e.g., sigmoid). Shown in Figure~\ref{fig2}, the gated self-attention (GSA) output is then written as
%   \begin{equation}
%     GSA(X)=SA(X)\odot G(X),
%   \end{equation}
% where $\odot$ denotes element-wise multiplication. This gating mechanism enables fine-grained modulation of attention responses at the token level.

% In addition to prefix conditioning, each transformer block is modulated by track-wise global conditioning embeddings derived from CLAP features. Let
%   \begin{equation}
%       C_i=[\,c_i^{(m)},c_i^{(u)},c_i^{(s)}\,]
%   \end{equation}
% denote the global conditioning embeddings for the mixture, sub-mixture, and source tracks. A lightweight conditioning network $\Psi(\cdot)$ maps $C_i$ to adaptive scale, shift, and gate parameters for the self-attention and feed-forward branches:
%   \begin{equation}
%     [s_{\mathrm{self}},b_{\mathrm{self}},g_{\mathrm{self}},s_{\mathrm{ff}},b_{\mathrm{ff}},g_{\mathrm{ff}}]=\Psi(C_i).
%   \end{equation}
% Following AdaLN-Zero~\cite{peebles2023scalable}, these parameters modulate each transformer block, providing text semantic control.
\subsection{Gated Self-Attention Modulation in DiT}
\label{4.3}
Our denoising backbone is a latent diffusion transformer \cite{evans2024long} that operates on the concatenated three-track latent sequence $X_i$. The self-attention operation is defined as
\begin{equation}
      Attn(Q,K,V)=\mathrm{softmax}\!\left(\frac{QK^\top}{\sqrt{d_k}}\right)V,
\end{equation}
where $d_k$ is the head dimension. To modulate attention, especially the attention sink phenomenon on timestep token, we employ a query-dependent output gate. Given the attention input $X$, a learnable gating function produces a head-specific element-wise gate
\begin{equation}
      G_i(X)=\sigma\!\left(XW^{G}_i\right),
\end{equation}
where $W^{G}_i$ refers to the learnable parameters of the $i$-th head and $\sigma$ is the sigmoid function. The gated attention output for each head is
\begin{equation}
      head_i=
      Attn(XW_i^Q,XW_i^K,XW_i^V) \odot G_i(X),
\end{equation}
where $\odot$ denotes element-wise multiplication. Shown in Figure~\ref{fig2}, the gated self-attention (GSA) output is then written as
\begin{equation}
      GSA(X)=
      \mathrm{Concat}\!\left(head_1,\ldots,head_h\right)W^{O}.
\end{equation}
This gating mechanism enables fine-grained modulation of attention responses at the token level before head fusion and output projection.

In addition to prefix conditioning, each transformer block is modulated by track-wise global conditioning embeddings derived from CLAP features. Let
\begin{equation}
        C_i=[\,c_i^{(m)},c_i^{(u)},c_i^{(s)}\,]
\end{equation}
denote the global conditioning embeddings for the mixture, sub-mixture, and source tracks. A lightweight conditioning network $\Psi(\cdot)$ maps $C_i$ to adaptive scale, shift, and gate parameters for the self-attention and feed-forward branches:
\begin{equation} [s_{\mathrm{self}},b_{\mathrm{self}},g_{\mathrm{self}},s_{\mathrm{ff}},b_{\mathrm{ff}},g_{\mathrm{ff}}]=\Psi(C_i).
\end{equation}
Following AdaLN-Zero~\cite{peebles2023scalable}, these parameters modulate each transformer block, providing text semantic control.

\subsection{Conditional Sampling via CFG}
We further improve robustness to incomplete conditions through a classifier-free conditioning strategy \cite{ho2022classifier}. During training, the global semantic conditioning $C_i$ is randomly dropped with probability $p_{1}$, while the quality conditioning $Q_i$ is occasionally replaced with a null condition with probability $p_{2}$. Note that when $Q_i$ is nullified by CFG dropout, the auxiliary quality mask (with probability $p_m$) is not applied to that sample, as masking an already-null token provides no additional supervisory signal. As a result, the model is exposed to four conditioning combinations, namely $(C_i, Q_i)$, $(C_i,\varnothing)$, $(\varnothing, Q_i)$, and $(\varnothing, \varnothing)$. In practice, these combinations are induced through condition removal rather than explicitly constructed as fully symmetric guidance branches.

Let $v_c = v_\theta(z_\tau, \tau \mid C_i, Q_i)$ denote the full conditional prediction, $v_t = v_\theta(z_\tau, \tau \mid C_i, \varnothing)$ the text-only prediction, $v_q = v_\theta(z_\tau, \tau \mid \varnothing, Q_i)$ the quality-only prediction, and $v_u = v_\theta(z_\tau, \tau \mid \varnothing, \varnothing)$ the unconditional prediction. The CFG prediction with independent guidance scales is:
  \begin{equation}
  v_\theta^{\text{CFG}}(z_\tau, \tau) =
  v_u + s_t \, (v_t - v_u) + s_q \, (v_q - v_u),
  \end{equation}
where $s_t$ and $s_q$ control the strength of text and quality guidance. Setting $(s_t, s_q) = (s, 0)$, $(0, s)$, $(0, 0)$, or $(s, s)$ yields text-only, quality-only, no guidance, or joint guidance respectively.

In the total generation task, the model samples the latent triplet under the joint condition $(C_i, Q_i)$:
  \begin{equation}\label{eq:mix_gen}
      \hat{z}^{(m)}, \hat{z}^{(u)}, \hat{z}^{(s)}
      \sim p_{\theta}\bigl(z^{(m)}, z^{(u)}, z^{(s)} \mid C_i, Q_i\bigr).
  \end{equation}
The synthesized mixture waveform is then obtained by decoding the generated mixture latent $\hat{z}^{(m)}$ through the
   pretrained VAE decoder $\mathit{D}$:
  \begin{equation}
      \hat{x}^{(m)} = D\!\left(\hat{z}^{(m)}\right).
  \end{equation}
In the source extraction task, the mixed latent representation $z^{(m)}$ is treated as known, and the target source is sampled under the joint condition $(C_i, Q_i)$:
  \begin{equation}\label{eq:source_extraction}
    \hat{z}^{(u)}, \hat{z}^{(s)} \sim
    p_{\theta}\!\left(z^{(u)}, z^{(s)} \mid z^{(m)}, C_i, Q_i\right),
    \quad
    \hat{x}^{(s)} = D\!\left(\hat{z}^{(s)}\right).
  \end{equation}
In the partial generation task, given the observed sub-mixture latent representation $z^{(u)}_0$, the model generates the missing source under the joint condition $(C_i, Q_i)$:
  \begin{equation}\label{eq:source_completion}
    \hat{z}^{(m)}, \hat{z}^{(s)}
    \sim
    p_{\theta}\!\left(z^{(m)}, z^{(s)} \mid z^{(u)}_{0}, C_i, Q_i\right).
  \end{equation}

\begin{algorithm}[t]
\caption{Quality-Aware Multi-Task DiT Joint Training}
\label{alg:quality_dit_training_fixed}
\begin{algorithmic}[1]
% \STATE Triplet latents $z_0 = \{z^{(m)}, z^{(u)}, z^{(s)}\} \in \mathbb{R}^{3C \times L}$, COCOLA score $q$, Text embedding $C$
% \STATE Denoising backbone $f_\theta$ with GSA weights $W^G$, condition net $\Psi(\cdot)$, aux head $g(\cdot)$
% \STATE Hyperparameters: $\gamma$, $\lambda_{\mathrm{q}}$, mask prob $p_m$, CFG prob $p_1, p_2$, drop prob $p_{\mathrm{drop}}$
\STATE \textbf{Input:} Triplet latents $z_0 = \{z^{(m)}, z^{(u)}, z^{(s)}\} \in \mathbb{R}^{3C \times L}$, COCOLA score $q$, text embedding $C$
\STATE \textbf{Output:} Training loss $\mathcal{L}$
\STATE \textbf{Model:} Denoising backbone $f_\theta$ with GSA weights $W^G$, condition net $\Psi(\cdot)$, aux head $g(\cdot)$
\STATE \textbf{Hyper:} $\gamma$, $\lambda_{\mathrm{q}}$, mask prob $p_m$, CFG prob $p_1, p_2$, drop prob $p_{\mathrm{drop}}$

\item[] \textbf{Quality Tokenization \& Masking Strategy:}
\STATE $\tilde{q} = \mathrm{clip}\big( (q - q_{\min}) / (q_{\max} - q_{\min}), \, 0, \, 1 \big)$
\STATE $\Gamma(\tilde{q}) = \big[ \sin(2\pi 2^f \tilde{q}), \cos(2\pi 2^f \tilde{q}) \big]_{f=0}^{F-1}$
\STATE $Q = \gamma \cdot \mathrm{MLP}_{\theta}(\Gamma(\tilde{q})) \in \mathbb{R}^{M \times D}$
\STATE $\mathbb{I}_{\mathrm{mask}} = \mathbb{I}(u \sim \text{Uniform}(0, 1) < p_m)$
\IF{$\mathbb{I}_{\mathrm{mask}} == 1$}
    \STATE $Q = [\mathbf{MASK}]_{M \times D}$
\ENDIF

\item[] \textbf{CFG Condition Dropping:}
\IF{$u_1 \sim \text{Uniform}(0, 1) < p_1$} 
    \STATE $C = \varnothing$ 
\ENDIF
\IF{$u_2 \sim \text{Uniform}(0, 1) < p_2$} 
    \STATE $Q \leftarrow Q_{\mathrm{null}}$, \quad $\mathbb{I}_{\mathrm{mask}} \leftarrow 0$ \hfill $\triangleright$ Learnable null token; disable aux loss
\ENDIF

\item[] \textbf{Multi-Task Timestep \& Noise Injection:}
\STATE Sample $\tau \sim \mathcal{U}([\tau_{\min},1])$, \quad $\epsilon \sim \mathcal{N}(0, \mathbf{I})_{3C \times L}$
\STATE Draw $r \sim \text{Uniform}(0, 1)$
\STATE $\boldsymbol{\tau} = (\tau, \tau, \tau)$ \textbf{if} $r > p_{\mathrm{drop}}$ \textbf{else} uniformly sample from $\{(0,\tau,\tau), (\tau,0,\tau), (\tau,\tau,0)\}$
\STATE $z_{\boldsymbol{\tau}} = \cos(\frac{\pi}{2}\boldsymbol{\tau}) \odot z_0 + \sin(\frac{\pi}{2}\boldsymbol{\tau}) \odot \epsilon$
\STATE $v_{\boldsymbol{\tau}} = \cos(\frac{\pi}{2}\boldsymbol{\tau}) \odot \epsilon - \sin(\frac{\pi}{2}\boldsymbol{\tau}) \odot z_0$

\item[] \textbf{Forward Propagation:}
\STATE $X = \big[\, \mathrm{Embed}(\boldsymbol{\tau}); \, Q; \, \mathrm{Linear}(z_{\boldsymbol{\tau}}) \,\big] \in \mathbb{R}^{(1 + M + N) \times D}$
% \STATE $\mathbf{G} = \sigma(X W^G) \in \mathbb{R}^{(1+M+N) \times H \times d_h}$
\STATE $\mathbf{G}_h = \sigma(X \, W^G_h) \in \mathbb{R}^{(1+M+N) \times d_h}, \quad h = 1,\ldots,H$ \hfill $\triangleright$ Head-specific gate, $W^G_h \in \mathbb{R}^{D \times d_h}$
\STATE $Y = f_\theta(X \mid \Psi(C); \mathbf{G})$ \hfill $\triangleright$ Forward modulation via GSA
\STATE $\hat{v}_{\boldsymbol{\tau}} = \mathrm{LinearOut}(Y[1+M : 1+M+N]) \in \mathbb{R}^{3C \times L}$
\STATE $\bar{h}^{(q)} = \mathrm{mean}(Y[1 : 1+M])$, \quad $\hat{q} = \sigma\big(g(\bar{h}^{(q)})\big)$

\item[] \textbf{Loss Computation:}
\STATE Split $\hat{v}_{\boldsymbol{\tau}}$ and $v_{\boldsymbol{\tau}}$ into tracks $\{m, u, s\}$ along channel dimension
\STATE $\mathcal{L}_{\mathrm{diff}} = \sum_{v \in \{m,u,s\}} \mathbb{I}(\boldsymbol{\tau}^{(v)} > 0) \cdot \|\hat{v}_{\boldsymbol{\tau}}^{(v)} - v_{\boldsymbol{\tau}}^{(v)}\|_2^2$
\STATE $\mathcal{L} = \mathcal{L}_{\mathrm{diff}} + \lambda_{\mathrm{q}} \cdot \mathbb{I}_{\mathrm{mask}} \cdot \|\hat{q} - \tilde{q}\|_2^2$
\RETURN $\mathcal{L}$
\end{algorithmic}
\end{algorithm}

\section{Experimental}
\subsection{Baselines}
We used the multi-track diffusion model MSDM \cite{mariani2023multi} and some of the most recent models, such as MSD-LD \cite{karchkhadze2025simultaneous} and MGE-LDM \cite{chae2025mge}, as benchmarks. These benchmarks are configured to support the synchronized modeling of standard stems (bass, drums, guitar, and piano) across total generation, partial generation, and source extraction tasks. 

To ensure a rigorous and fair comparison, all baseline metrics are strictly evaluated on our designated test splits under aligned inference conditions, following the evaluation paradigms established in  MGE-LDM \cite{chae2025mge} that all evaluations are rigorously aligned to identical audio specifications, utilizing 12.8-second clips in a 16 kHz audio:
\begin{itemize}
    \item MSDM: We utilize its official pretrained checkpoint and implementation. To match its 22 kHz processing rate, we upsample our 16 kHz evaluation clips for inference and downsample the generated outputs back to 16 kHz.
    \item MSD-LD: As no checkpoint is publicly released, we retrained it using the official codebase under identical configurations.
    \item MGE-LDM: We utilize its officially released codebase and checkpoint, evaluated on 12.8-second audios at 16 kHz.
\end{itemize}

\subsection{Datasets }
Our training and evaluation utilize Slakh2100 \cite{manilow2019cutting}, Musdb18 \cite{MUSDB18HQ} and MoisesDB \cite{pereira2023MoisesDB}. Slakh2100 is a dataset of multi-track audio and aligned MIDI for music source separation and multi-instrument automatic transcription. Musdb18 is a dataset of 150 full lengths music tracks of different genres along with their isolated drums, bass, vocals and other stems. MoisesDB consists of 240 tracks from 45 artists across 12 musical genres. For each song, we used all stems and organized them into triples. Processed audio combinations are encoded into latent representations prior to model training to reduce computational overhead. By evaluating all models on the exact same unseen evaluation splits of Slakh2100, Musdb18, and MoisesDB, we mitigate potential performance discrepancies stemming from mismatched test distributions, thereby guaranteeing a fair architectural benchmark.

\subsection{Implementation Details }
% Our model is founded on Stable Audio \cite{evans2025stable}, which includes an autoencoder and a diffusion model based on DiT \cite{peebles2023scalable} and MEG-LDM \cite{chae2025mge}. For the composite condition of each triplet, the quality condition applies the method outlined in Chapter IV, while the text condition vector employs text embeddings derived from the pre-trained CLAP model (HTSAT-base) on the music-audioset-epoch-15-esc-90.14 dataset. The GSA module processes quality scores through QT-Encoder, a sinusoidal positional encoding followed by an MLP that generates learnable quality tokens. The model can be trained on one or multiple NVIDIA RTX A6000 GPUs.
Our model is built upon Stable Audio \cite{evans2025stable} and MGE-LDM \cite{chae2025mge}, leveraging a VAE alongside a DiT backbone \cite{peebles2023scalable} for effective audio generation. The VAE encodes 16 kHz audio, yielding latents of dimension $C=64$ and sequence length $T=100$ frames. For multi-track modeling, the three tracks are each encoded independently and concatenated along the channel dimension, producing a joint input of $3C=192$ to the DiT. This compression discards redundant phase information while retaining semantic and acoustic structure, enabling the 24-layer/48-head DiT to model cross-track dependencies efficiently. For the composite condition of each triplet, the quality condition applies the method outlined in Chapter IV, while the text condition vector employs text embeddings derived from the pre-trained CLAP model (HTSAT-base) on the music-audioset-epoch-15-esc-90.14 dataset. The GSA module processes quality scores through QT-Encoder, a sinusoidal positional encoding followed by an MLP that generates learnable quality tokens. The model can be trained on one or multiple NVIDIA RTX A6000 GPUs.

\begin{table*}[t]
\centering
\caption{Comprehensive comparison across total generation, partial generation, and source extraction tasks. Total generation results measured by FAD (lower is better). Partial generation results measured by COCOLA score (higher is better). and FAD (lower is better). Source extraction results measured by  Log-Mel L1 distance (lower is better). “Random” means for each audio track, one stem is randomly selected as the source stem, and the rest are selected as submix
stems}
\resizebox{\textwidth}{!}{
\begin{tabular}{l c ccccc ccccc ccccc}
\toprule

& \multicolumn{1}{c}{Total} 
& \multicolumn{5}{c}{Partial (COCOLA $\uparrow$)} 
& \multicolumn{5}{c}{Partial (FAD $\downarrow$)} 
& \multicolumn{5}{c}{Extraction (Log-Mel L1 $\downarrow$)} \\

\cmidrule(lr){2-2}
\cmidrule(lr){3-7}
\cmidrule(lr){8-12}
\cmidrule(lr){13-17}

Model 
& FAD 
& bass & drum & guitar & piano & random
& bass & drum & guitar & piano & random
& bass & drum & guitar & piano & random\\

\midrule

MSDM 
& 4.2123 
& 44.06 & 44.20 & 42.92 & 44.26 & 43.86
& 1.19 & 0.79 & 0.91 & 1.35 & 1.06
& 5.26 & 2.14 & 7.08 & 5.86 & 5.09\\

MSD-LD 
& 1.3086 
& 48.46 & 47.58 & 46.58 & 49.92 & 48.14
& \textbf{0.33} & 0.34 & \textbf{0.49} & 1.08 & \textbf{0.60}
& 2.92 & 4.02 & 4.08 & 4.05 & 3.77\\

MGE-LDM 
& 0.4723 
& 47.31 & 55.32 & 41.90 & 52.95 & 49.37
&1.14 &1.50 &3.75 &2.47  & 2.22
&1.67 &0.83 &\textbf{2.61} &1.77  & 1.60\\

Ours 
& \textbf{0.3728} 
& \textbf{48.68} & \textbf{58.10} & \textbf{52.66} & \textbf{55.66} & \textbf{53.88}
& 0.54 & \textbf{0.31} & 1.13 & 1.15 & 0.73
& \textbf{1.07} & \textbf{0.58} & 2.72 & \textbf{1.20} & \textbf{1.26}\\

\bottomrule
\end{tabular}
}
\label{tab:combined_results}
\end{table*}
Training is divided into two stages: first, training the autoencoder, then training the latent diffusion transformer. The frozen COCOLA \cite{ciranni2025cocola} module uses pre-trained COCOLA-HP-v1. The model is trained for 200K steps with a batch size of 64. We use AdamW optimizer with learning rate of $5\times10^{-5}$, beta coefficients $(0.9, 0.999)$, and weight decay of $0.001$. The learning rate follows an inverse square root decay schedule (InverseLR) with inverse gamma $\gamma=10^6$, power 0.5, and warmup factor 0.99. All models utilized 12.8-second, 16 kHz audio clips. The latent diffusion transformer consists of 24 layers with a per-track embedding dimension of 512 (joint dimension $D=1536$ for three tracks), 48 attention heads, and GroupNorm normalization, adopting v-prediction as the diffusion objective. The timestep sampler uses a scrambled Sobol quasi-random sequence over $[\tau_{\min}, 1]$ for low-discrepancy coverage and a minimum timestep $\epsilon=0.02$. Gradient clipping with a maximum norm of 1.0 and an exponential moving average (EMA) of model weights is applied during training. Then, in the processes of sampling and patching, our CFG parameter configurations comprised a guidance scale of 2.0 and a dropout probability $p_1,p_2=0.4,0.4$ for each track. Diffusion samples were generated with DDIM sampling and 250 inference steps.The quality encoder consists of a two-layer MLP with hidden dimension $d_{\mathit{hidden}}=256$, LayerNorm, and GELU activation, mapping the $2F$-dimensional sinusoidal encoding to the transformer hidden dimension $D=3 \times 512$. The quality tokens are generated via sinusoidal positional embedding with a learnable scale $\gamma$ initialized to 1.0, with $p_m = 0.15$ during training and an auxiliary quality loss weight of 0.1. During inference, we support three guidance modes: joint guidance $(s_t, s_q) = (2.0, 2.0)$, text-only guidance $(s_t, s_q) = (2.0, 0.0)$, and quality-only guidance $(s_t, s_q) = (0.0, 2.0)$. 

Experimental results are based on the downstream tasks of total generation, partial generation, and source extraction. Both total and partial generation utilize the FAD score computed via VGGish embeddings \cite{hershey2017cnn} and the FAD metric \cite{kilgour2018fr}. A lower FAD score indicates that the generated audio more closely matches the distribution of authentic audio. For the partial generation task, the COCOLA score was additionally employed. 
% Calculations using COCOLA's computation pipeline to assess the similarity between the model-generated track $x$ and the provided reference tracks $x^{m}-x^{s}$, where the track type in $x$ matches the track type in $x^{s}$, are equal to the track type in $x^{s}$. 
The COCOLA score evaluates the musical coherence between the model-generated track $x$ and the provided accompaniment tracks $x^{s}$. Higher values indicate greater coherence and coordination between the two tracks. Note that COCOLA is used solely as an evaluation metric for partial generation, while FAD provides an independent complementary assessment, ensuring no circularity with the quality token training objective. For source separation tasks, SI-SDR\cite{le2019sdr} is not applicable to our model, as the VAE decoder introduces reconstruction differences in both phase and waveform that preclude direct sample-level comparison with the reference signal. We use the Log-Mel L1 distance mentioned in MSD-LD \cite{karchkhadze2025simultaneous} for evaluation. Smaller values indicate closer proximity to the reference. The implementation is available at \url{https://github.com/locacaca/QAG-LDM}.


\subsection{Comparison with advanced models}
The FAD between the mixed audio from the test set and the resulting audio track is calculated. Table \ref{tab:combined_results} demonstrates that our model attains superior performance on this criterion, slightly outperforming MGE-LDM and significantly outperforming MSD-LD and MSDM. This shows that in total generation tasks, models can generate multi-track audio mixes that closely resemble the distribution of grounded music. Notably, compared to the strongest baseline MGE-LDM (0.4723), ours achieves a $21.07\%$ relative reduction in FAD score, demonstrating a substantial leap in rendering higher musicality and audio quality.

In the partial generation task, models generate tracks featuring bass, drums, guitar, piano, and random stems. Regarding the FAD metric, MSDM shows higher FAD values across all instrument generation tasks, while our method performs best when generating drums. This difference may be due to the drums' rhythm being stronger than other stems, which aligns well with human auditory perception. Notably, MSD-LD exhibits a slight numerical lead in FAD scores for certain instrument tracks during partial generation tasks. This could be attributed to the fact that the FAD metric primarily evaluates the global statistical alignment between generated samples and real-world distributions, whereas ours prioritizes track-to-track coherence and localized feature refinement. As evidenced by COCOLA scores in Table \ref{tab:combined_results}, compared with MGE-LDM, our model brings a $9.14\%$ relative improvement in COCOLA score, indicate that our model offers a more robust framework for maintaining musical logic and inter-track coordination.

We also evaluated the source extraction task.Ours yields a Log-Mel L1 distance of 1.26 on the random task. This represents a $21.25\%$ improvement over the competitive MGE-LDM and a $75.25\%$ substantial reduction over MSDM. Our method shows better performance in the separation of both guitar and bass. From the model architecture, we can see that latent diffusion-based models, including ours and MGE-LDM, yield better results in the source extraction task. And traditional diffusion models, including MSDM, demonstrate lower Log-Mel L1 metrics for all stems. The only exception occurs in the guitar track, where ours shows a marginal degradation in terms of the Log-Mel L1 distance compared with MEG-LDE. This localized trade-off can be fundamentally attributed to the highly volatile acoustic properties of guitar signals and our model's global optimization priorities. Our quality-aware constraint imposes a minor smoothing penalty on the low-level spectral distance of the highly dynamic guitar stem, yet yields a substantially higher perceived acoustic fidelity and robust multi-track integration on the whole.

\begin{table*}[t!]
  \centering
\caption{Ablation study on quality control architectures. Q1 denotes joint embedding of text and quality information, including "SED", "KL", and "linear" interpolation-based variants. Q2 denotes gated dual-branch cross-attention. Q3 (ours) denotes prefix-based quality control. For Q3, we further vary the number of quality tokens, the token position in the prefix sequence, the quality encoding method, and the auxiliary quality regression loss. "Added Param" indicates added parameters. Detailed architectural differences are described in the text.}
  \begin{tabular}{lccc|cccccc}
  \toprule
  Variant & Arch. & M & Aux loss & Added Param &COCOLA$\uparrow$ & FAD(p)$\downarrow$ &
  Log-Mel L1$\downarrow$  & FAD(t)$\downarrow$ & MSE loss$\downarrow$\\
  \midrule
  baseline & - & -- & $\times$  & - & 49.37 & 2.22 & 1.60 & 0.47 & 0.304\\ 
  Q1-SED & Q1 & -- & $\times$  & 4.7 & 47.12 & 2.45 & 1.68 & 0.51 & 0.315\\
  Q1-KL  & Q1 & -- & $\times$  & 4.7 & 48.05 & 2.38 & 1.65 & 0.49 & 0.348\\
  Q1-LI  & Q1 & -- & $\times$  & 1.6 & 49.10 & 2.25 & 1.61 & 0.48 & 0.305\\
  Q2     & Q2 & -- & $\times$  & 37.8 & 50.44 & 1.85 & 1.82 & 0.42 & 0.282\\
  Q3-1   & Q3 & 1 & $\times$  & 5.9 & 49.23 & 1.45 & 1.60 & 0.39 & 0.312\\
  Q3-2   & Q3 & 1 & $\checkmark$  & 6.7 & \textbf{52.12} & \textbf{0.98} & \textbf{1.61} & \textbf{0.39} & 0.228\\
  Q3-3   & Q3 & 2 & $\checkmark$  & 13.8 & 51.24 & 1.15 & 1.85 & 0.41 & \textbf{0.226}\\
  \bottomrule
  \end{tabular}
  \label{tab:quality_architecture_ablation}
\end{table*}

\subsection{Ablation Experiment}
To examine the influence of each design decision on generation performance, this section conducts comprehensive ablation experiments. For all tasks except total generation, we employ a randomized sampling strategy in which a subset of stems is selected as $x^s$.

\subsubsection{Quality Control Architectures and Parameters}
We experimented with various quality control architectures, including directly integrating quality information with text information (Q1), conducting cross-attention between quality information and audio (Q2), and tokenizing quality information as prefix tokens (Q3, ours). 

For Q1, we fuse the text condition embedding $e_t$ and the quality condition embedding $e_q$ into a unified representation $e_c$ by solving $e_c=\arg\min_{e}\sum_{k\in\{\mathrm{t},\mathrm{q}\}}\lambda_k D_\phi(e,e^k)$, where $D_\phi(\cdot,\cdot)$ denotes a Bregman divergence induced by a strictly convex function $\phi(\cdot)$, and $\lambda_k$ controls the relative importance of each condition, thereby encouraging the joint embedding to remain semantically consistent with both text and quality information. And we also explore multiple variants of the joint embedding formulation by instantiating $D_\phi$ with different divergence measures, including the squared Euclidean distance (SED), KL divergence, and a simple linear interpolation $e_{c}=\alpha e_{t}+(1-\alpha)e_{q}$, thereby examining how different geometric assumptions affect the fusion of text and quality conditions.

For Q2, we adopt a dual-branch cross-attention mechanism where the audio latent $\mathit{X}$ attends independently to the text embedding $e_t$ and the quality embedding $e_q$ via two cross-attention modules, and their outputs are fused as $\tilde{\mathit{X}}=\mathit{X}+\sigma(1-\mathit{g}_{\mathit{self}})\odot\mathit{Attn}(\mathit{adaLN}(\mathit{X}))$, $\mathit{F}=\tilde{\mathit{X}}+[\mathit{CrossAttn}(\tilde{\mathit{X}},e_t)+\mathit{CrossAttn}(\tilde{\mathit{X}},e_q)]\mathit{W}^O$.

For Q3, we additionally perform ablations within the proposed prefix-based quality control architecture to identify the most effective configuration. In particular, we study the number of quality tokens, the position of the quality tokens in the prefix sequence, the quality encoding method, and the use of the auxiliary quality regression loss. These experiments help disentangle whether the gains of Q3 come from the prefix-token formulation itself, from the way the scalar quality score is encoded, or from the additional supervisory signal, and therefore provide a more fine-grained understanding of the proposed quality control mechanism.

\begin{table}[t!]
    \centering
    \caption{Comparison of model metrics under different quality fraction settings, where COCOLA and FAD(p) represent the COCOLA score and Fréchet audio distance metric for partial generation tasks, respectively; Log-Mel L1 denotes the Log-Mel L1 distance between the separated audio track and the target track in the separation task, and FAD(t) indicates the Fréchet audio distance metric for the full generation task.}
    \begin{tabular}{lcccc}
        \toprule
        $\tilde{q}_i$  & COCOLA↑ & FAD(p)↓ & Log-Mel L1↓ & FAD(t)↓ \\
        \midrule
        0.2     & 47.09   & 1.0169 & 1.5372 & 0.4002 \\
        0.5     & 52.89   & 0.9383 & 1.4890 & 0.3861 \\
        0.8     & \textbf{53.87}  & \textbf{0.7278} & \textbf{1.2615} & \textbf{0.3728} \\
        \bottomrule
    \end{tabular}
    \label{tab:q}
\end{table}
\begin{table}[t!]
    \centering
    \caption{Performance comparison under different guidance scales for text and quality conditions.}
    \begin{tabular}{lcccc}
        \toprule
        Mode & COCOLA↑ & FAD(p)↓ & Log-Mel L1↓ & FAD(t)↓ \\
        \midrule
        $s_t = 1.0$, $s_q = 0.0$     & 45.21   & 0.78 & 1.71 & 0.49 \\
        $s_t = 2.0$, $s_q = 0.0$     & 45.93   & 0.71 & 1.41 & 0.45 \\
        $s_t = 0.0$, $s_q = 1.0$     & -   & - & - & 0.82 \\
        $s_t = 0.0$, $s_q = 2.0$     & -   & - & - & 0.72 \\
        $s_t = 1.0$, $s_q = 2.0$     & 53.12   & 0.79 & 1.48 & \textbf{0.37} \\
        $s_t = 1.0$, $s_q = 1.0$     & 51.89   & \textbf{0.70} & 1.39 & 0.40 \\
        $s_t = 2.0$, $s_q = 1.0$     & 51.42   & 0.73 & \textbf{1.21} & 0.39 \\
        $s_t = 2.0$, $s_q = 2.0$     & \textbf{53.87}  & 0.73 & 1.26 & \textbf{0.37} \\
        \bottomrule
    \end{tabular}
    \label{tab:mode_ablation}
\end{table}

The results are shown in Table \ref{tab:quality_architecture_ablation}. The additional parameters required for the method that simply integrates the quality conditions and text conditions into the training and reasoning process are relatively few (Q1), but the improvement in performance on the standard MSDM task is very small, and the MSE loss remains at a high level. Adding an additional cross-attention method can improve the performance of some models, but the parameter quantity has increased by 70\% (Q2). The method of encoding the quality conditions into prefix tokens of length M (Q3-1, Q3-2, Q3-3) adds relatively fewer additional parameters, 
and we can observe that the aux loss has a significant impact on the semantic consistency, bringing $5.87\%$ relative improvement in COCOLA score (from 49.23 of Q3-1 to 52.12 of Q3-2) and $32.41\%$ relative improvement in FAD(p) (from 1.45 of Q3-1 to 0.98 of Q3-2), which proves the effectiveness of the masked-quality supervision strategy
Increasing M instead will weaken the model's performance (Q3-3). We speculate that this is because the quality semantics already have an appropriate semantic information density in the encoding format of length 1, and increasing the length will weaken the expression ability of the quality tokens.

\begin{figure*}[t]
    \centering

    % 第一行
    \begin{minipage}[b]{0.48\textwidth}
        \centering
        \includegraphics[width=\textwidth]{layer_18_ungated_heatmap.png}
    \end{minipage}
    \hfill
    \begin{minipage}[b]{0.48\textwidth}
        \centering
        \includegraphics[width=\textwidth]{layer_21_ungated_heatmap.png}
    \end{minipage}

    \vspace{2pt} % 控制上下间距（可调）

    % 第二行
    \begin{minipage}[b]{0.48\textwidth}
        \centering
        \includegraphics[width=\textwidth]{layer_18_gated_heatmap.png}
    \end{minipage}
    \hfill
    \begin{minipage}[b]{0.48\textwidth}
        \centering
        \includegraphics[width=\textwidth]{layer_21_gated_heatmap.png}
    \end{minipage}

    \caption{Average attention map weights for each head (the upper is the baseline; the lower is ours). We selected the most representative layer for the heat map comparison. Layer 21 in the baseline model demonstrates a strong attention sink on the timestep token, which is substantially reduced by the gate.}
    \label{fig:heatmap_compare}
\end{figure*}

We also aim to confirm the guiding role of quality conditions and explore how different quality values affect the model's performance across tasks. We selected the representative quality parameters based on Figure \ref{fig:cocola_distribution}. The lower peak corresponds to regions with poorer data coordination, while the higher peak indicates areas with superior coordination. Thus, we selected normalized values of 0.2, 0.5, and 0.8 to represent the low, median, and high peak scenarios, respectively, to assess model performance. Table \ref{tab:q} shows that on the semantic consistency metric COCOLA, 
% model performance continuously improves with increasing quality values, and the distribution also approaches the mean more closely. Regarding audio quality, both FAD(p) and Log-Mel L1 achieve optimal results at $\tilde{q}_i = 0.8$, 
model performance continuously improves with increasing quality values. Notably, when $\tilde{q}_i$ increases from 0.2 to 0.8, the COCOLA score exhibits a $14.40\%$ relative growth. Regarding audio quality, both FAD(p) and Log-Mel L1 achieve optimal results at $\tilde{q}_i = 0.8$, where the Log-Mel L1 distance is reduced by $17.93\%$,
indicating that higher quality condition parameters help generate samples closer to the distribution of real audio. This monotonic improvement across semantic and multi-source distribution metrics demonstrates that the model exhibits a controlled sensitivity to $\tilde{q}$, verifying that the quality tokens successfully guide the prior learning process rather than being ignored. Conversely, FAD(t) remains relatively stable across the three quality values, which highlights the model's inherent resilience in terms of temporal feature consistency. This implies that while fine-grained semantic harmony and acoustic fidelity are sensitive to $\tilde{q}$ modulation, the fundamental structural and temporal framework of the generated music remains robust against varying quality conditions. Overall, $\tilde{q}_i = 0.8$ yields the best performance across all settings, indicating superior semantic consistency and harmony in its generated content. The fact that performance safely scales with higher $\tilde{q}$ without catastrophic degradation or overfitting further confirms the selection resilience of our quality-aware conditioning framework. Therefore, in subsequent sections, we use this value for comparisons with other models, and metric calculations follow the above description.

\begin{table*}[t!]
      \centering
      \caption{Ablation study on gating architectures. "+Q"
  indicates that quality-aware prefix conditioning and the auxiliary
  quality regression loss are enabled. "Gate Pos." denotes the
  position where the gate is inserted in the self-attention pathway.
  "Type" indicates whether the gate is shared across heads or
  head-specific. T-Attn denotes the mean attention mass assigned to
  the timestep token in self-attention.}
      \begin{tabular}{lccc|ccccccc}
          \toprule
Variant & +Q & Gate Pos. & Type & Added Param &COCOLA$\uparrow$ & FAD(p)$\downarrow$ & Log-Mel L1$\downarrow$ & FAD(t)$\downarrow$ & T-Attn$\downarrow$ & MSE loss$\downarrow$ \\
          \midrule
          G1 & $\times$ & None & -- & - & 49.37 & 2.22 & 1.60 & 0.47 & 0.33 & 0.3041 \\
          G1 & $\checkmark$ & None & -- & 7 & 52.12 & 0.98 & 1.51 & 0.39 & 0.35& 0.2750 \\
          G2 & $\times$ & Pre-Attn & Shared & 124 & 49.88 & 1.85 & 1.65 & 0.45 & 0.22 & 0.2985 \\
          G2 & $\checkmark$ & Pre-Attn & Shared & 131 & 51.56 & 0.87 & 1.42 & 0.39 & 0.20 & 0.2488 \\
          G3 & $\times$ & Post-Proj & Shared & 124 & 50.12 & 1.72 & 1.50 & 0.44 & 0.28 & 0.2910 \\
          G3 & $\checkmark$ & Post-Proj & Shared & 131 & 52.34 & 0.84 & 1.20 & 0.38 & 0.26 & 0.2465 \\
          G4-S & $\times$ & Pre-Proj & Shared & 124 & 50.45 & 1.75 & 1.44 & 0.43 & 0.20 & 0.2882 \\
          G4-S & $\checkmark$ & Pre-Proj & Shared & 131 & 52.88 & 0.79& \textbf{1.18} & \textbf{0.37} & 0.17 & 0.2442 \\
          G4-H & $\times$ & Pre-Proj & Head-Specific & 446 & 51.10 & 1.58 & 1.45 & 0.42 & 0.11 & 0.2640 \\
          G4-H & $\checkmark$ & Pre-Proj & Head-Specific & 453 & \textbf{53.87} & \textbf{0.72} & 1.26 & \textbf{0.37} & \textbf{0.09} & \textbf{0.2426} \\
          \bottomrule
      \end{tabular}
      \label{tab:gating_ablation}
\end{table*}

\subsubsection{Conditioning Mode}
We conduct ablation experiments on different conditioning combinations during inference. We evaluated various modes. Table \ref{tab:mode_ablation} presents the results. The text-only mode (setting $s_t > 0, s_q = 0$) leverages CLAP embeddings to control text semantics, while the quality-only mode (setting $s_t = 0, s_q > 0$) relies solely on the GSA branches with quality tokens. The joint mode (setting $s_t > 0, s_q > 0$) combines both pathways for compound control.

When the quality guidance scale is zero ($s_q = 0.0$), the model relies on text condition. Comparing rows with $s_q = 0.0$, we observe that increasing $s_t$ improves COCOLA scores, etc., confirming that text guidance enhances semantic alignment. However, inference solely based on text condition does not perform optimally in all situations.

When the text guidance scale is zero ($s_t = 0.0$), the model depends exclusively on GSA branches with quality tokens. However, in this situation, the model theoretically loses the meanings of the partial generation task and the source extraction task. This is because when the text guidance scale is 0, the model is unable to specify the correct stems as a necessary input condition for the task.

The joint guidance mode with $s_t = 2.0$ and $s_q = 2.0$ achieves the best performance across most metrics, with a COCOLA score of 53.87 and FAD(t) of 0.37. This validates that decoupling text and quality into independent pathways enables flexible and complementary control. Furthermore, intermediate combinations (e.g., $s_t = 2.0$, $s_q = 1.0$) provide a smooth trade-off between semantic alignment and audio quality, allowing users to balance these aspects according to their specific requirements (FAD(p) and Log-Mel L1 perform best). The above experimental results show that our quality control and text quality control have a complementary relationship. And during inference, enhancing the CFG parameter of one method will not reduce the performance of the other control method.

\begin{figure}[t]
    \centering
    % 如果需要裁剪，取消注释下行并调整数值：trim = left bottom right top
    % \includegraphics[trim=0cm 0cm 0cm 0cm, clip, width=\linewidth]{cocola.pdf}
    \includegraphics[width=\linewidth]{attention_sink_curve_ungated_epoch120_vs_attention_sink_curve_gated_epoch120_lineplot.png}
    \caption{Timestep token attention score. The experimental results show that the gated mechanism exhibits extremely strong denoising capabilities in the intermediate layers (with a sink reduction of over 95$\%$ in layer 18), successfully releasing the semantic modeling space of the model. While in the output layer, the model retains a proportion (with a sink reduction of 36$\%$ in layer 23) of attention on the timestep token. The average attention of the timestep token for all layers decreased by 73$\%$. }
    \label{fig:attention_sink_epoch}
\end{figure}

\subsubsection{Ablation on Gating Designs and Quality Conditioning}

During the experiment, we found that the attention score of the timestep token $T_i$ at the first position of $X_i$ was excessively high. The baseline-defined timestep is crucial for distinguishing the task, but excessive timestep token attention in each layer would cause the model to pay less attention to the quality latents and audio latents. We attempted to alleviate this problem by using a gating mechanism. To identify the most effective gating design in our gated latent diffusion transformer, we conduct an ablation study comparing variants that differ in three aspects: 1) whether gating is used, 2) where the gate is inserted in the self-attention computation, and 3) whether the gate is shared across attention heads or predicted independently for each head. We additionally report each variant with and without quality-aware conditioning, denoted by "+Q." Here, "+Q" indicates that the model uses the quality prefix tokens together with the auxiliary quality regression loss, whereas variants without "+Q" remove the quality-control pathway entirely. The results are summarized in Table \ref{tab:gating_ablation}.

No gating (G1) serves as the baseline. Pre-attention gating (G2) applies the gate to the query representation $X_i$ before the scaled dot-product attention is computed. Post-projection gating (G3) applies the gate after the multi-head attention output has already been projected back to the model dimension. Pre-projection gating (G4) applies the gate directly to the multi-head attention output before the final output projection. This is the design adopted in our final model.

Performance is reported in Table~\ref{tab:gating_ablation}. The gating mechanisms added at various positions (G2, G3, G4) have different impacts on performance. Overall, all variants outperform the baseline (G1), and G4 performs the best. And the architectures that have enabled quality-aware prefix conditioning and the auxiliary
quality regression loss perform better than those that have not been enabled. G2 modulates information too early in the attention pipeline, which may undesirably constrain the interactions required for accurate denoising and source reconstruction. This variant yields the COCOLA score 51.56, FAD(t) 0.39, FAD(p) 0.98, and Log-Mel L1 1.92. We can observe that the improvement is not significant. Although G3 preserves the original attention computation itself, the gate acts on a representation that has already been mixed across heads, which weakens fine-grained control over head-wise responses. By gating the attention response before head fusion, G4 can more directly suppress irrelevant responses and preserve informative ones during denoising.

We also analyze the timestep token attention score, denoted as T-Attn, to examine the attention sink behavior on the timestep token of each variant. Higher T-Attn values indicate that the model assigns excessive attention mass to the initial token, which is undesirable for stable sequence modeling \cite{xiao2023efficient}. MSE loss represents the optimization objective of the model for the reconstruction error of the main task. As reported in Table \ref{tab:gating_ablation}, the ungated baseline exhibits a stronger attention sink on timestep token tendency with 0.33, while the gated variants reduce this value respectively. 

After determining the optimal gate position, we further compare two parameterizations of this design: a shared gate, which applies the same gating pattern across all heads, and a head-specific gate, which predicts separate gates for different heads. In particular, the head-specific pre-projection gate yields the best performance with the lowest T-Attn and MSE loss, suggesting that fine-grained gating helps the model bypass irrelevant token interactions and allocate more capacity to salient audio-semantic structures. When combined with quality control (+Q), this effect is further reflected in improved generation fidelity and semantic consistency.

From the layer-by-layer comparison results in Figure \ref{fig:attention_sink_epoch}, it can be observed that after introducing the gated mechanism, the average attention sink of the timestep token has significantly decreased (0.33 vs. 0.09), indicating that this design effectively alleviates the problem where the model highly relies on the timestep token to stabilize the attention distribution. In the intermediate layers, the sharp reduction in sink values means that the model successfully decentralized the attention mechanism, transferring computing resources from meaningless placeholders to sequence features with actual semantics, thereby strengthening the feature transformation and alignment capabilities. As it approaches the output layer, although the sink value has decreased, it still remains at a relatively high level (0.56 vs. 0.43). This reflects that the model has evolved a task-adaptive characteristic: in deep networks, the model actively converts time step information into a "dynamic bias" to ensure that the generated results strictly follow the global constraints of time steps in the final output stage. This "intermediate layer decoupling and output layer anchoring" distribution pattern strongly proves that the gated mechanism not only improves the model's representation efficiency but also does not compromise its perception of key control signals.

\section{Conclusion and Discussion}
We propose QAG-LDM, a versatile framework that incorporates quality-aware conditioning into the multi-track diffusion paradigm through a gated attention mechanism. By introducing the QT-Encoder, we successfully integrated a learnable quality signal into the latent space, enabling the model to transition from a quality-agnostic training paradigm to quality-aware acoustic distributions. Experimental results across multiple benchmarks demonstrate that our framework not only achieves state-of-the-art performance in total generation and source extraction tasks but also significantly enhances the coherence of partial generation. This confirms that explicitly modeling the intrinsic quality relationships within structured data is crucial for aligning generative outputs with human auditory perception.

A key contribution of our work is the identification and mitigation of the attention sink phenomenon on timestep tokens within DiT. Our investigation revealed that the excessive concentration of attention on the timestep token often limits the model's ability to capture fine-grained semantic features. Through the implementation of the GSA, we provided an adaptive modulation pathway that redistributes focus across the latent sequence. This architectural refinement, coupled with the auxiliary quality-prediction loss, ensures a more stable and balanced learning process. 

While QAG-LDM effectively manages global quality signals, future research could explore more granular, multi-dimensional quality control, such as separate conditioning for rhythmic precision and timbral purity. Additionally, expanding the framework to support long-form music generation while maintaining the efficiency of the gating mechanism remains a critical challenge. By refining the synergy between quality-aware representation and efficient structural modeling, we aim to bridge the gap between AI-driven content creation and professional-grade music production, ultimately providing more intuitive and robust tools for the creative community.

\bibliographystyle{IEEEtran}
\bibliography{taslp} % 对应你的 ref.bib
\end{document}